\documentclass[11pt, a4paper, twoside]{article}

% --- Packages ---
\usepackage[utf8]{inputenc}
\usepackage{amsmath, amssymb, amsthm}
\usepackage{geometry}
\usepackage{tcolorbox}
\usepackage{tikz-cd}
\usepackage{fancyhdr}
\usepackage{booktabs}

% --- Page Layout & Styling ---
\geometry{
	top=1in, 
	bottom=1in, 
	left=1.25in, 
	right=1.25in
}

\pagestyle{fancy}
\fancyhf{}
\fancyhead[LE,RO]{\thepage}
\fancyhead[LO]{\small \textit{Lecture Notes: Advanced Coding Theory}}
\fancyhead[RE]{\small \textit{Classification and Galois Cohomology}}
\renewcommand{\headrulewidth}{0.4pt}

% Custom boxes
\newtcolorbox{conceptbox}[1][]{
	colback=blue!5!white, colframe=blue!75!black,
	fonttitle=\bfseries, title=#1, arc=3pt, boxrule=1pt,
	left=10pt, right=10pt, top=10pt, bottom=10pt
}
\newtcolorbox{theorybox}[1][]{
	colback=gray!5!white, colframe=gray!75!black,
	fonttitle=\bfseries, title=#1, arc=3pt, boxrule=1pt,
	left=10pt, right=10pt, top=10pt, bottom=10pt
}

% --- Macros ---
\newcommand{\F}{\mathbb{F}}
\newcommand{\C}{\mathcal{C}}
\newcommand{\Aut}{\text{Aut}}
\newcommand{\Gal}{\text{Gal}}

\begin{document}
	
	\begin{center}
		{\Large \textbf{The Classification Problem and Galois Cohomology}} \\[0.5em]
		{\normalsize \textit{Measuring the Twists of Equivalent Codes}}
	\end{center}
	\vspace{1em}
	
	A major problem in classical coding theory is \textbf{classification}: knowing when two seemingly different codes are actually the exact same code in disguise. To solve this, mathematicians use Galois Cohomology to rigorously classify all possible ``twists'' of a code over a finite field.
	
	\section*{1. The Definition of Equivalence}
	
	Two linear codes $C_1, C_2 \subseteq \F_q^n$ are not considered strictly identical unless their generator matrices are row-equivalent. However, they are \textbf{isometric} (or equivalent) if one can be transformed into the other simply by rearranging the transmission channel or scaling the symbols.
	
	\begin{conceptbox}[Monomial Equivalence]
		Two codes are equivalent if there exists an $n \times n$ \textbf{monomial matrix} $M$ (a permutation matrix multiplied by a diagonal matrix of non-zero scalars) such that:
		$$ C_2 = C_1 M $$
		This means $C_1$ and $C_2$ have the exact same length, dimension, and minimum distance. They represent the same geometric subspace, just viewed through a different coordinate basis. 
	\end{conceptbox}
	
	To classify codes, we must define the \textbf{Automorphism Group} $\Aut(C)$, which consists of all monomial matrices $M$ that map the code $C$ back onto itself ($CM = C$). The size and structure of $\Aut(C)$ dictate how symmetric the code is.
	
	\section*{2. The Principle of Galois Descent}
	
	Often, it is easy to classify codes over a large, algebraically closed field (or a large extension field $\F_{q^m}$), but very difficult to classify them over the smaller base field $\F_q$. 
	
	Suppose we have a code $C$ defined over $\F_{q^m}$. We want to know: \textit{How many distinct codes exist over the base field $\F_q$ that become equivalent to $C$ when we extend our scalars up to $\F_{q^m}$?}
	
	
	
	This is the problem of \textbf{Galois Descent}. We are looking for ``forms'' or ``twists'' of $C$. If two codes over $\F_q$ become isomorphic over $\F_{q^m}$, their difference is purely a structural twist governed by the Galois group $\Gal(\F_{q^m} / \F_q)$.
	
	\section*{3. The First Galois Cohomology Set ($H^1$)}
	
	\begin{theorybox}[Classifying the Twists]
		Let $L/K$ be a Galois extension of fields (e.g., $\F_{q^m} / \F_q$). Let $C$ be a linear code over $K$, and let $\Aut(C_L)$ be its automorphism group over the extension field $L$.
		
		The Galois group $\Gamma = \Gal(L/K)$ acts continuously on $\Aut(C_L)$. We define a \textbf{1-cocycle} as a continuous map $c: \Gamma \to \Aut(C_L)$ satisfying the cocycle condition:
		$$ c_{\sigma \tau} = c_\sigma \cdot \sigma(c_\tau) \quad \text{for all } \sigma, \tau \in \Gamma $$
		
		Two cocycles $c$ and $c'$ are cohomologous if there exists some $a \in \Aut(C_L)$ such that $c'_\sigma = a^{-1} c_\sigma \sigma(a)$. 
		The set of equivalence classes of these cocycles forms the \textbf{First Galois Cohomology Set}, denoted:
		$$ H^1(\Gal(L/K), \Aut(C_L)) $$
	\end{theorybox}
	
	\section*{4. The Coding Theoretical Dictionary}
	
	What does this abstract cohomological set actually compute? It serves as a strict counting mechanism for classification.
	
	\begin{itemize}
		\item \textbf{The Base Object ($C$):} A known, standard linear code over $\F_q$.
		\item \textbf{The Automorphisms ($\Aut(C_L)$):} The symmetries of the code over the larger field. If the code is highly symmetric, it can be twisted in many ways. If it is rigid, it cannot be twisted.
		\item \textbf{The Cocycle ($c_\sigma$):} A specific ``recipe'' for twisting the coordinate basis of the code as you descend from the extension field down to the base field.
		\item \textbf{The Cohomology Set ($H^1$):} There is a rigorous bijection between the elements of $H^1(\Gal(L/K), \Aut(C_L))$ and the isomorphism classes of codes over $\F_q$ that become equivalent to $C$ over $L$.
	\end{itemize}
	
	\textbf{Conclusion:} If $H^1$ contains 5 elements, it means there are exactly 5 distinct, non-equivalent codes over your hardware's base field $\F_q$ that all share the exact same geometric structure when viewed from the extension field. Cohomology guarantees you have found every possible classification without duplicates.
	
\end{document}